\section{\href{https://doi.org/10.1103/PhysRevA.84.061404}{Time delay between photoemission from the 2p and 2s subshells of neon}}

\subsection*{Supervisor: Prof. Hugo van der Hart}


The R-matrix incorporating time (RMT) method is a method developed recently 
for solving the time-dependent Schroedinger equation for multielectron atomic 
systems exposed to intense short-pulse laser light. We have employed the RMT 
method to investigate the time delay in the photoemission of an electron 
liberated from a 2p orbital in a neon atom with respect to one released from 
a 2s orbital following absorption of an attosecond xuv pulse. Time delays due 
to xuv pulses in the range 76--105 eV are presented. For an xuv pulse at the 
experimentally relevant energy of 105.2 eV, we calculate the time delay to 
be 10.2$\pm$1.3 attoseconds (as), somewhat larger than estimated by other 
theoretical calculations, but still a factor of 2 smaller than experiment. 
We repeated the calculation for a photon energy of 89.8 eV with a larger 
basis set capable of modeling correlated-electron dynamics within the neon 
atom and the residual Ne ion. A time delay of 14.5$\pm$1.5 as was observed, 
compared to a 16.7$\pm$1.5 as result using a single-configuration 
representation of the residual Ne\textsuperscript{+} ion. 
